\section{白盒：Dijkstra 与 Bellman-Ford}\label{sec:white}

\subsection{记号}

全文讨论无向连通图 $G=(V,E,w)$，$|V|=n$，$|E|=m$，边权 $w:E\to\mathbb{R}_{>0}$。
记 $d_G(s,t)$ 为 $s$ 到 $t$ 的最短路长度，$\mathrm{dist}_G(s,\cdot)$ 为单源距离向量。
对 $s\in V$，记 $N(v)$ 为 $v$ 的邻居集合。

\subsection{Dijkstra 的正确性与复杂度}

Dijkstra 维护一个“已确定”集合 $S$，每轮取出 $d$ 值最小的未确定点加入 $S$，并松弛它的出边。
其正确性依赖一条归纳不变式：\emph{每次加入 $S$ 的点的 $d$ 值就是最终最短距离}。
证明的关键是边权非负——若存在负权，已经确定的点可能被后来的路径更新，不变式失效。
这也解释了为什么实现里必须显式拒绝负权：修复前我们实测过，负权输入下 Dijkstra 会静默返回源点距离为负、父指针成环的结果（回溯函数于是死循环）；现在 \texttt{dijkstra} 入口会直接抛 \texttt{ValueError}。
Dijkstra 与 Bellman-Ford 的原始文献见~\cite{dijkstra1959,bellman1958}，标准的正确性证明可参考~\cite{cormen2009clrs}。

用二叉堆实现时，Dijkstra 的复杂度是 $O((n+m)\log n)$；在 $m\propto n$ 的稀疏图上即 $O(n\log n)$。
实测（图~\ref{fig:complexity}）在 $n=800$ 时为 7.4\,ms（同规模 Bellman-Ford 1.70\,s）。

\subsection{Bellman-Ford 的 DP 形式}\label{sec:bf}

Bellman-Ford 的常见写法是“松弛 $n-1$ 轮”。但\textbf{第 $k$ 轮到底算出了什么}，往往被含糊带过。它的准确语义是：

\begin{proposition}[轮数语义]\label{prop:rounds}
设 $d^{(0)}(s)=0$，$d^{(0)}(v)=+\infty\ (v\ne s)$，并按
\begin{equation}
d^{(k+1)}(v)=\min\Bigl\{d^{(k)}(v),\ \min_{(u,v)\in E}\bigl[d^{(k)}(u)+w(u,v)\bigr]\Bigr\}
\label{eq:bf}
\end{equation}
\emph{同步}更新（第 $k+1$ 轮只读第 $k$ 轮的值），则 $d^{(k)}(v)$ 恰为“从 $s$ 到 $v$、至多 $k$ 条边”的最短路长度。
\end{proposition}

\begin{proof}
对 $k$ 归纳。$k=0$ 时只有 $s$ 自身可达且需要 0 条边，成立。
设对 $k$ 成立。任取一条从 $s$ 到 $v$、至多 $k+1$ 条边的最短路 $P$。
若 $P$ 至多 $k$ 条边，则由归纳假设 $d^{(k)}(v)$ 已不大于 $|P|$，而式~\eqref{eq:bf} 取 $\min$ 后不会变大；
若 $P$ 恰有 $k+1$ 条边，设其最后一条边为 $(u,v)$，则 $P$ 的前缀是至多 $k$ 条边的 $s$--$u$ 路，由归纳假设 $d^{(k)}(u)\le |P|-w(u,v)$，于是 $\min$ 项给出 $d^{(k+1)}(v)\le|P|$。
另一方面，式~\eqref{eq:bf} 的每个候选值都对应一条“至多 $k+1$ 条边”的走道，故 $d^{(k+1)}(v)$ 不会小于真正的最短路。
\end{proof}

\begin{figure}[H]
\centering
\includegraphics[width=0.46\textwidth]{complexity.png}
\caption{三类精确算法的实测复杂度（ER 图，平均度 6，每个规模取 3 张图的中位数）。双对数拟合斜率：Dijkstra 1.09（理论 $n\log n$ 在该区间约 1.18）、Bellman-Ford 2.12（理论 $O(nm)$ 即 2）、Floyd 2.96（理论 $O(n^3)$ 即 3）。计时随机器负载波动，此处取提交的 \texttt{results/complexity.csv}。}
\label{fig:complexity}
\end{figure}

命题~\ref{prop:rounds} 看起来只是教科书里的一个注脚，但它是本文的关键：因为它把“轮数”与“边数”绑定在了一起，而消息传递的“层数”与“跳数”绑定在一起。
两者一旦对上，第~\ref{sec:black}~节的等价就是自然的推论，而不是类比。

\subsection{一个最小的验证}

为了确认实现没有把“同步”写成“原地”，我们用 DFS 枚举所有简单路径作为独立真相源，在 100 张 $n\le 7$ 的随机图上逐轮比对。
结果 100/100 一致：第 $k$ 轮确实等于“至多 $k$ 条边”的最短路。若改成原地更新（同一轮内就传播多条边），这条性质会立刻失效。

\subsection{复杂度实测的诚实说明}

图~\ref{fig:complexity} 的拟合斜率与渐近阶同量级但不完全吻合：Dijkstra 1.09 低于理论值 1.18，而 Bellman-Ford 2.12 略高于 2，Floyd 2.96 略低于 3。
原因不是实现错了，而是规模不够大：每次调用的固定开销（入参校验、堆分配、数组预分配）在 $n\le800$ 时占比可观，拟合出来的斜率测的是常数项而不是阶。
我们用“最大两点之间的局部指数”做了补充：Dijkstra 1.10、Bellman-Ford 2.05、Floyd 3.57。
报告里同时给出拟合斜率、局部指数和原始时间表，读者可以自己判断。


